% document formatting
\documentclass[10pt]{article}
\usepackage[utf8]{inputenc}
\usepackage[left=1in,right=1in,top=1in,bottom=1in]{geometry}
\usepackage[T1]{fontenc}
\usepackage{xcolor}
\usepackage[table,xcdraw]{xcolor}

% math symbols, etc.
\usepackage{amsmath, amsfonts, amssymb, amsthm}

% lists
\usepackage{enumerate}
\usepackage{enumitem}
\usepackage{tabularx}
\usepackage{multirow}

% images
\usepackage{graphicx} % for images
\usepackage{tikz}
\usepackage[dvipsnames]{xcolor}

% code blocks
\usepackage{minted, listings} 

% verbatim greek
\usepackage{alphabeta}
\DeclareMathOperator*{\argmin}{arg\,min}
\DeclareMathOperator*{\argmax}{arg\,max}

\newcommand{\dd}{\text{d}}


\graphicspath{{./assets/images/Week 12}}

\title{03-621 Week 12 \\ \large{Advanced Quantative Genetics}}
 
\author{Aidan Jan}

\date{\today}

\begin{document}
\maketitle

\section*{Epigenetics and Epigenomics}
\subsection*{Parts of Genes}
\subsubsection*{Promoters}
\begin{itemize}
	\item Major on-off switches in genes
	\item Recruit transcription initiation complexes and RNA polymerase
	\item Most elements are upstream of the start site
	\item Some promoters lack these core elements
\end{itemize}
\subsubsection*{Enhancers}
\begin{itemize}
	\item Sequences that interact with activator proteins to boost transcription
	\item Frequently tissue-, stage-, or condition-specific
	\item Usually <20bp long
	\item Usually upstream of the promoter, but can be \textbf{far} upstream, downstream, or within a gene (in introns)
\end{itemize}
\subsubsection*{Insulator Elements}
\begin{itemize}
	\item Long-range looping can occur due to tissue-specific enhancers interacting long range
	\item Insulator elements prevent enhancer-promoter interactions
	\item Can restrict the activity of enhancers to defined chromatin domains (don't activate nearby genes).
\end{itemize}
\subsubsection*{Locus Control Regions}
Can affect the transcription of multiple, distantly located genes, ``choosing between them''
\begin{itemize}
	\item For example, in the $\beta$-globin gene cluster, the same LCR can associate with different $\beta$-globin genes at different times in development, activating them sequentially.
\end{itemize}

\subsection*{Transcriptional Condensates: A New Regulatory Layer}
\begin{itemize}
	\item Transcriptional condensates are phase-separated clusters of transcription factors, mediator, and RNA Pol II
	\item Driven by intrinsically disordered regions (IDRs) and multivalent interactions
	\item Responsible for concentrating transcription machinery at active loci
	\item Enhancers nucleate condensates to regulate gene expression
	\item Highly dynamic: assmeble/disassemble in response to signals
\end{itemize}

\subsection*{Chromatin}
\begin{itemize}
	\item Chromatin compacts the DNA so it fits and remains organized in the nucleus. 
	\item This structure also makes it hard for transcription factors and transcriptional machinery to access DNA
\end{itemize}
\begin{center} 
	\includegraphics*[width=0.7\textwidth]{W12_1.png} 
\end{center}

\section*{Epigenetics}
\begin{itemize}
	\item Regulation of Gene Expression at the level of chromatin structure, histone modification and DNA methylation
	\item ON/OFF epigenetic states can be established in response to:
	\begin{itemize}
	    \item Developmental program (cell type differentiation)
	    \item Environmental factors
    \end{itemize}
    \item Epigenetic regulatory states (ON/OFF) can be inherited through somatic cell divisions
    \begin{itemize}
	    \item Provides a ``cellular memory'' of previous decisions in the cell lineage
    \end{itemize}
\end{itemize}
\textbf{Chromatin states arise from DNA and chromatin modifications.}
\begin{center} 
	\includegraphics*[width=0.7\textwidth]{W12_2.png} 
\end{center}
Different types of histone methylation are associated with open and condensed chromatin.\\\\
\textbf{Contributions by DNA and Histone modification to different chromatin states}
\begin{center} 
	\includegraphics*[width=0.7\textwidth]{W12_3.png} 
\end{center}
Histones H3 and H4 undergo many modifications on their positively charged N-terminal tails.  H3 Lysines K9 and K27 can be:
\begin{itemize}
	\item Methylated (inactive genes)
	\item Acetylated (active genes)
\end{itemize}
\begin{center} 
	\includegraphics*[width=0.9\textwidth]{W12_4.png} 
\end{center}
\subsubsection*{Examples of Histone Modifications Characteristic of different Chromatin States}
\begin{center}
\begin{tabular}{|c|ccccc|cc|}
\hline
\rowcolor[HTML]{CBCEFB} \cellcolor[HTML]{CBCEFB} & \multicolumn{5}{c|}{\cellcolor[HTML]{CBCEFB}\textbf{Euchromatin}} & \multicolumn{2}{c|}{\cellcolor[HTML]{CBCEFB}\textbf{Heterochromatin}} \\ \cline{2-8} 
\rowcolor[HTML]{EAEBFF} \cellcolor[HTML]{CBCEFB} & \multicolumn{2}{c|}{\cellcolor[HTML]{EAEBFF}Promoters} & \multicolumn{2}{c|}{\cellcolor[HTML]{EAEBFF}Enhancers} & Gene Bodies & \multicolumn{1}{c|}{\cellcolor[HTML]{EAEBFF}} & \cellcolor[HTML]{EAEBFF} \\ \cline{2-6}
\rowcolor[HTML]{EAEBFF} \cellcolor[HTML]{CBCEFB}\multirow{-3}{*}{\begin{tabular}[c]{@{}c@{}}\cellcolor[HTML]{CBCEFB}\textbf{Amino}\\ \cellcolor[HTML]{CBCEFB}\textbf{Acid}\end{tabular}} & \multicolumn{1}{c|}{\cellcolor[HTML]{EAEBFF}Active} & \multicolumn{1}{c|}{\cellcolor[HTML]{EAEBFF}Inactive} & \multicolumn{1}{c|}{\cellcolor[HTML]{EAEBFF}Active} & \multicolumn{1}{c|}{\cellcolor[HTML]{EAEBFF}Inactive} & Inactive & \multicolumn{1}{c|}{\cellcolor[HTML]{EAEBFF}\multirow{-2}{*}{Facultative}} & \multirow{-2}{*}{\cellcolor[HTML]{EAEBFF}Constitutive} \\ \hline
H3K4 & \multicolumn{1}{c|}{\begin{tabular}[c]{@{}c@{}}H3K4me2,\\ H3K4me3\end{tabular}} & \multicolumn{1}{c|}{} & \multicolumn{1}{c|}{\begin{tabular}[c]{@{}c@{}}H3K4me1,\\ H3K4me2\end{tabular}} & \multicolumn{1}{c|}{} &  & \multicolumn{1}{c|}{} &  \\ \hline
H3K9 & \multicolumn{1}{c|}{H3K9ac} & \multicolumn{1}{c|}{H3K9me3} & \multicolumn{1}{c|}{} & \multicolumn{1}{c|}{\begin{tabular}[c]{@{}c@{}}H3K9me2,\\ H3K9me3\end{tabular}} & \begin{tabular}[c]{@{}c@{}}H3K9me2,\\ H3K9me3\end{tabular} & \multicolumn{1}{c|}{H3K9me2} & H3K9me3 \\ \hline
H3K27 & \multicolumn{1}{c|}{H3K27ac} & \multicolumn{1}{c|}{H3K27me3} & \multicolumn{1}{c|}{H3K27ac} & \multicolumn{1}{c|}{} &  & \multicolumn{1}{c|}{H3K27me3} &  \\ \hline
H4K12 & \multicolumn{1}{c|}{H4K12ac} & \multicolumn{1}{c|}{} & \multicolumn{1}{c|}{} & \multicolumn{1}{c|}{} &  & \multicolumn{1}{c|}{H4K12ac} & H4K12ac \\ \hline
H4K20 & \multicolumn{1}{c|}{} & \multicolumn{1}{c|}{} & \multicolumn{1}{c|}{} & \multicolumn{1}{c|}{} &  & \multicolumn{1}{c|}{} & H4K20me3 \\ \hline
\end{tabular}
\end{center}

\subsubsection*{CpG density and DNA methylation levels across a typical gene}
\begin{center} 
	[insert 15]
\end{center}
\textbf{CpG Islands:}
\begin{itemize}
	\item 300-3000 nucleotides long, with high G/C content.
	\item Methylation of C leads to silencing of gene expression
	\item Often found in transposons (suppresses hopping).
	\item 60-70\% of CpGs are methylated in human somatic cells.
\end{itemize}

\section*{Epigenetics}
\begin{itemize}
	\item Heritable change in something ``in addition to'' or ``on top of'' DNA sequence.
	\item Methylation of C's in CpG islands by DNA Methyltransferase (DNMT)
	\item Results in local silencing of transcription
	\item Methylation is inherited through somatic cell divisions, through the action of ``maintenance'' methylases
\end{itemize}
\begin{center} 
	[insert 16]
\end{center}

\subsection*{DNA Methylation and Demethylation at CpG in Mammalian Cells}
\begin{enumerate}
	\item \textbf{De novo methylation} can establish a repressed state in response to developmental program or environmental factors
	\item Methylation can be lost through subsequent rounds of cell division
	\item \textbf{Maintenance methylation} maintains repressed state through somatic cell divisions.  Accounts for cellular ``memory'' of past regulatory decisions
\end{enumerate}
\begin{center} 
	[insert 19]
\end{center}

\subsection*{DNA Methylation in Neanderthals}
We suspect that neanderthals diverged from humans through epigenetic evolution.
\begin{itemize}
	\item Neanderthal genome sequenced in 2010
	\item 5-methyl-cytosine naturally deaminates to thymidine.
	\item Methylation sites ``look'' like mutations, compared to modern human genome sequences.
	\item The C $\rightarrow$ T ``mutation'' shows the methylation status in Neanderthals
	\item Over 99\% of genome shows no methylation different from present-day humans
\end{itemize}
Comparing the HOXD genes in DNA of osteoblasts (bone cells)
\begin{itemize}
	\item Neanderthals had \textit{more} methylation at HOXD9 and HOXD10 genes.
	\item This may have led to \textit{less} expression, accounting for differences in bone structure.
	\item LOF mutations in these genes lead to similar differences in mouse, human
\end{itemize}

\subsection*{Methylation Paradox}
If methylation silences gene expression, and it is heritable through cell divisions, how can new progeny develop?\\\\
Answer: Methylation status is controlled dynamically
\begin{itemize}
	\item Erased during development of germ cells
	\item Reset (where appropriate) during subsequent differentiation of gametes, fertilization, and development
	\item The repressed state is \textbf{reversible} in response to the developmental program (e.g., during gamete differentiation and early embryonic development, or in response to environmental factors)
\end{itemize}

\subsection*{Changes in DNA Methylation During Mammalian Development}
\begin{center} 
	[insert 26]
\end{center}
\begin{itemize}
	\item CpG islands are demethylated in PGCs and in preimplantation embryos
	\item CpG islands become methylated during development for differentiation (as appropriate for sperm, oocytes, embryo, and adult somatic cells).
	\item Deficient DNA methylation is associated with Neural Tube Defects (folic acid vitamin supplement increases methylation - reduces risk)
\end{itemize}

\subsection*{Examples of Regulation by DNA Methylation During Development}
\begin{center} 
    \begin{tabular}{|c|c|c|}
    \hline
    \multirow{4}{*}{\textbf{\begin{tabular}[c]{@{}c@{}}Independent of \\ parent of origin\end{tabular}}} & \textbf{\begin{tabular}[c]{@{}c@{}}X-inactivation in\\ somatic cells\end{tabular}} & \begin{tabular}[c]{@{}c@{}}Random inactivation of most genes on either the \\ paternal or maternal X\end{tabular} \\ \cline{2-3} 
    & \textbf{\begin{tabular}[c]{@{}c@{}}Production of\\ cell-specific\\ Ig and T-cell receptors\end{tabular}} & \begin{tabular}[c]{@{}c@{}}Each mature B or T cell makes Ig or T-cell receptor\\ chains using only one allele.  Once a functional chain\\ is made by gene rearrangement at one randomly \\ selected allele, a feedback mechanism inhibits\\ further rearrangements\end{tabular} \\ \cline{2-3} 
    & \textbf{\begin{tabular}[c]{@{}c@{}}Production of\\ cell-specific olfactory\\ receptors\end{tabular}} & \begin{tabular}[c]{@{}c@{}}Each olfactory neuron expresses a single allele of just\\ a single olfactory receptor (OR) gene (selected from\\ several hundred OR genes) so that it fires in\\ response to one specific odorant only.  Depends on\\ competition for a single monoallelic enhancer.\end{tabular} \\ \cline{2-3} 
    & \textbf{\begin{tabular}[c]{@{}c@{}}Stochastic and\\ environmentally induced\end{tabular}} & \begin{tabular}[c]{@{}c@{}}May be quite common?  Intense interest in context \\ of cancer, obesity, and diabetes.\end{tabular} \\ \hline
    \multirow{2}{*}{\textbf{\begin{tabular}[c]{@{}c@{}}Dependent on\\ parent of origin\end{tabular}}} & \textbf{\begin{tabular}[c]{@{}c@{}}X-inactivation in\\ placenta\end{tabular}} & Paternal X is always inactivated \\ \cline{2-3} 
    & \textbf{Genomic imprinting} & \begin{tabular}[c]{@{}c@{}}Some genes are expressed only from the paternally\\ inherited homolog and some only from the maternally\\ inherited homolog.  Examples: IGF2, H19, 15q11, etc.\end{tabular} \\ \hline
    \end{tabular}
\end{center}

\subsection*{Imprinting Occurs During Germ Cell Differentiation}
\begin{itemize}
	\item In mammals, we would not expect trans-generational inheritance of epigenetic changes acquired in response to environmental factors.
	\item But some specific epigenetic marks are transferred through the sperm or egg $\rightarrow$ \textbf{Imprinting}: parent-of-origin silencing / expression
\end{itemize}

\subsection*{Imprinting}
\begin{itemize}
	\item Imprinted genes are methylated differentially in maternally and paternally-derived chromosomes.
	\item Silenced loci remain off during development.
	\item Exhibited by several hundred human genes.
	\item Therefore, we only have one \textbf{active} copy of those genes.  It is either the maternal or paternal copy, but it depends on the gene.
	\item The silenced allele can \textbf{switch} as it passes through a new germline.
\end{itemize}

\subsection*{Some Phenotypic Traits Exhibit Imprinted Transmission}
\begin{center} 
	[insert 33]
\end{center}
\begin{itemize}
	\item \textbf{Parent-of-origin effects:} the disease phenotype associated with a DNA sequence variant is only observed when that variant is inherited from the parent of a particular sex (in this case, the female parent).
\end{itemize}

\subsection*{Pathogenesis Can Result from Underexpression \textit{or} Overexpression of Imprinted Genes}
\begin{center} 
    \begin{tabular}{|lllll|}
    \hline
    \multicolumn{1}{|c|}{\multirow{2}{*}{\textbf{Syndrome}}} & \multicolumn{1}{c|}{\multirow{2}{*}{\textbf{Clinical Features}}} & \multicolumn{1}{c|}{\multirow{2}{*}{\textbf{\begin{tabular}[c]{@{}c@{}}Molecular Basis of \\ Imprinting Disorder\end{tabular}}}} & \multicolumn{2}{c|}{\textbf{Cause}} \\ \cline{4-5} 
    \multicolumn{1}{|c|}{} & \multicolumn{1}{c|}{} & \multicolumn{1}{c|}{} & \multicolumn{1}{c|}{\textbf{UPD$^*$}} & \multicolumn{1}{c|}{\textbf{Other}} \\ \hline
    \multicolumn{5}{|l|}{\textbf{Pathogenesis due to under-expression of imprinted genes}} \\ \hline
    \multicolumn{1}{|l|}{Prader-Willi} & \multicolumn{1}{l|}{\begin{tabular}[c]{@{}l@{}}Developmental delay\\ hypotonia, hyperphagia,\\ obesity, heart disease\end{tabular}} & \multicolumn{1}{l|}{\begin{tabular}[c]{@{}l@{}}Lack of active alleles \\ for imprinted snoRNA \\ genes at 15q11.2\end{tabular}} & \multicolumn{1}{l|}{\begin{tabular}[c]{@{}l@{}}Mat. chr 15\\ $\sim 25\%$\end{tabular}} & \begin{tabular}[c]{@{}l@{}}Del pat. 15q11-13  \\ $\sim 70\%$\end{tabular} \\ \hline
    \multicolumn{1}{|l|}{Angelman} & \multicolumn{1}{l|}{\begin{tabular}[c]{@{}l@{}}Developmental delay (severe)\\ no speech\\ epilepsy; ataxia\end{tabular}} & \multicolumn{1}{l|}{\begin{tabular}[c]{@{}l@{}}Lack of active allele for \\ imprinted UBE3A gene \\ at 15q11.2\end{tabular}} & \multicolumn{1}{l|}{\begin{tabular}[c]{@{}l@{}}Pat. chr 15\\ $\sim 25\%$\end{tabular}} & \begin{tabular}[c]{@{}l@{}}Del mat. 15q11-13 \\ $\sim 75\%$\end{tabular} \\ \hline
    \multicolumn{5}{|l|}{\textbf{Pathogenesis due to over-expression of imprinted genes}} \\ \hline
    \multicolumn{1}{|l|}{\begin{tabular}[c]{@{}l@{}}Beckwith-\\ Wiedeman\end{tabular}} & \multicolumn{1}{l|}{\begin{tabular}[c]{@{}l@{}}Macrosomia/overgrowth;\\ macroglossia\\ umbilical defect\end{tabular}} & \multicolumn{1}{l|}{\begin{tabular}[c]{@{}l@{}}Biallelic expression of IGF2 \\ at 11p15.5 (normally silenced \\ on maternal 11) Biallelic \\ expression of ncRNA that \\ silences CDKN1C growth \\ suppressor\end{tabular}} & \multicolumn{1}{l|}{Pat. 11} & \begin{tabular}[c]{@{}l@{}}Loss of maternal ICR2$^{**}$\\ methylation $\sim 50\%$\\ Gain of maternal ICR1$^{**}$\\ methylation $\sim 5\%$\\ CDKN1C mutation $\sim 5\%$\end{tabular} \\ \hline
    \end{tabular}
\end{center}




\end{document}